\chapter{Podsumowanie i~wnioski}

\section{Podsumowanie wyników zrealizowanych eksperymentów}
W~ramach niniejszej pracy magisterskiej przeprowadzono kompleksowe i~rygorystyczne metodologicznie badania empiryczne, których nadrzędnym celem była ilościowa i~jakościowa ocena wpływu 7~strategii przygotowania danych na skuteczność oraz stabilność czterech architektur klasyfikatorów tabelarycznych (Random Forest, XGBoost, Gaussian Naive Bayes oraz Perceptron Wielowarstwowy MLP).

Eksperymenty zrealizowano na 5 zróżnicowanych domenowo zbiorach danych (3 zbiorach medycznych: \textit{Zapalenia naczyń}, \textit{Choroby serca}, \textit{Cukrzyca Indian Pima}, zbiorze telekomunikacyjnym \textit{Rezygnacje klientów} oraz finansowym \textit{Ryzyko kredytowe}). Każdy ze zbiorów poddano kontrolowanemu procesowi degradacji jakości na 5 poziomach zanieczyszczeń ($10\%, 20\%, 30\%, 40\%, 50\%$), symulując jednoczesne występowanie trzech typów anomalii: losowych braków danych, ekstremalnych wartości odstających oraz błędnych kodów kategorii jakościowych. Ewaluację przeprowadzono w~procedurze powtarzanej stratyfikowanej walidacji krzyżowej, a~istotność różnic zweryfikowano testami rangowymi Friedmana i~Wilcoxona na poziomie istotności $\alpha = 0{,}05$.

\section{Główne wnioski z~przeprowadzonych eksperymentów}
Na podstawie szczegółowej analizy wyników eksperymentalnych, zestawień metryk $\overline{\mathrm{AUC}}$, analizy zysków $\Delta \mathrm{AUC}$ oraz wyników testów nieparametrycznych sformułowano następujące kluczowe wnioski:

\begin{enumerate}
    \item \textbf{Skuteczność potoków wieloetapowych przy narastającej degradacji danych}:  
    Brak jakichkolwiek procedur czyszczących (\textit{raw}) prowadzi do gwałtownego spadku jakości predykcji wszystkich modeli wraz ze wzrostem stopnia uszkodzenia danych. Zastosowanie pełnego potoku trójstopniowego (\textit{all} oraz \textit{all\_knn}) skutecznie amortyzuje ten spadek, pozwalając modelom drzewiastym (RF, XGBoost) zachować $\mathrm{AUC} > 0{,}80 - 0{,}86$ na zbiorach medycznych nawet przy $50\%$ uszkodzonych atrybutów. W~testach statystycznych Wilcoxona strategie te wykazały istotną statystycznie przewagę nad danymi surowymi w~$80\%$ badanych konfiguracji ($16$ na $20$ przypadków). Kluczowe znaczenie ma zachowanie kolejności operacji: wcześniejsza izolacja wartości odstających (konwersja do \texttt{NaN}) przed imputacją i~standaryzacją zapobiega zniekształceniu parametrów rozkładu.

    \item \textbf{Wyraźna asymetria wrażliwości badanych klasyfikatorów}:  
    Eksperymenty uwidoczniły silne zróżnicowanie reakcji modeli na zanieczyszczenia i~konieczność dopasowania potoku do architektury. Sieć neuronowa MLP na danych nieoczyszczonych przy $50\%$ uszkodzeń całkowicie traciła zdolność decyzyjną ($\mathrm{AUC} \approx 0{,}504$, co odpowiada losowemu zgadywaniu), natomiast wdrożenie potoku z~normalizacją (\textit{all} lub \textit{fill\_norm}) podnosiło jej skuteczność do poziomu $0{,}642$ (zysk $\Delta \mathrm{AUC} > +0{,}13 - +0{,}20$). Podobnie klasyfikator GaussianNB tracił zdolności dyskryminacyjne przez zawyżenie wariancji rozkładu normalnego przez outliery, zyskując średnio $+0{,}08 - +0{,}15$ AUC po ich odfiltrowaniu. Z~kolei modele zespołów drzew (Random Forest i~XGBoost) wykazały wysoką naturalną odporność, zyskując głównie na samej eliminacji błędów grubych i~poprawnej imputacji.

    \item \textbf{Porównanie metod imputacji (kompromis pomiędzy dokładnością a~czasem)}:  
    Porównanie strategii uzupełniania braków wykazało, że algorytm \texttt{KNNImputer} bez uprzedniej filtracji anomalii (\textit{fill\_knn}) ulega zaburzeniu metryki odległości euklidesowej. Dopiero w~połączeniu z~eliminacją outlierów (\textit{all\_knn}) metoda ta osiągnęła lekką przewagę nad medianą na zbiorach o~silnych korelacjach cech (\textit{Diabetes}, \textit{Serce}). Narzut czasowy obliczeń był jednak około 12-krotnie wyższy, co sprawia, że prosta imputacja medianą w~ramach potoku \textit{all} stanowi najbardziej optymalny i~szybki wybór w~praktycznych zastosowaniach.
\end{enumerate}

\section{Kierunki dalszego rozwoju badań}
Jako perspektywy kontynuacji i~rozwoju podjętej tematyki badawczej wskazuje się:
\begin{itemize}
    \item zbadanie skuteczności zaawansowanych metod imputacji wielokrotnej, takich jak algorytm MICE (ang. \textit{Multivariate Imputation by Chained Equations}) oraz lasy imputacyjne \textit{MissForest},
    \item ewaluację odporności nowoczesnych architektur głębokich dedykowanych danym tabelarycznym (m.in.~TabNet, FT-Transformer, SAINT),
    \item analizę wpływu zanieczyszczeń w~warunkach dryfu pojęciowego (ang. \textit{concept drift}) w~strumieniowym przetwarzaniu danych w~czasie rzeczywistym.
\end{itemize}
